\documentclass[a4paper,12pt]{article}
\usepackage{graphicx} % For inserting images
\usepackage[serbian]{babel} % Serbian language support
\usepackage[a4paper, margin=2.5cm]{geometry} % Improved page layout
\usepackage{fancyhdr} % Custom headers/footers
\usepackage{titlesec} % Better section titles
\usepackage{tocloft} % Better TOC spacing
\usepackage{biblatex} % Bibliography management
\usepackage{csquotes} % Recommended for biblatex
\usepackage[T1]{fontenc} % For đ

\setlength{\headheight}{14.5pt} % Fix fancyhdr warning

\pagestyle{fancy}
\fancyhf{} % Clear default header/footer
\fancyhead[L]{\small \leftmark} % Header left
\fancyfoot[C]{\thepage} % Page number in center footer

\titleformat{\section}{\large\bfseries}{\thesection}{1em}{}

\addbibresource{references.bib} % Reference file

\title{Privatnost u doba veštačke inteligencije}
\author{Marko Lazarević}
\date{Mart 2025}

\renewcommand{\contentsname}{Sadržaj}

\begin{document}

\maketitle
\newpage
\tableofcontents
\newpage

\section{Uvod}

Veštačka inteligencija (\textbf{AI}) je nešto sa čim se sve češće susrećemo – postala je deo savremenog života i mnogima predstavlja sastavni deo svakodnevice.

Evolucija veštačke inteligencije (\textbf{AI}) i mašinskog učenja (\textbf{ML}) predvodila je digitalnu transformaciju u poslednjoj deceniji. AI i ML su postigli značajne proboje, počevši od \textit{nadgledanog učenja}, a zatim su se brzo razvijali kroz \textit{nenadgledano}, \textit{polunadgledano}, \textit{učenje potkrepljivanjem} i \textit{duboko učenje}. Najnovija granica AI tehnologije stigla je u obliku \textit{generativne veštačke inteligencije}~\cite{10198233}.

\subsection{Izazovi i zabrinutosti u vezi sa AI}

Iako veštačka inteligencija donosi brojne koristi, postoje i mnoga \textbf{pravna} i \textbf{etička pitanja} koja se mogu postaviti kada se govori o njenoj upotrebi. Na primer, kada razgovaramo sa čet-botom (\textit{ChatGPT}) i podelimo s njim strategije i planove za budućnost svoje kompanije, postavlja se pitanje da li do tih informacija može doći konkurentska firma. Takođe, kada svoju fotografiju obradimo pomoću AI \textbf{alata za ulepšavanje}, nameće se sumnja da li ta slika može završiti u rukama osobe koja bi pokušala da nam ukrade identitet. Još jedno važno pitanje jeste kako se pravno reguliše razvoj, upotreba i odgovornost u vezi sa veštačkom inteligencijom – ko snosi odgovornost ako AI donese pogrešnu odluku i kako se štite prava korisnika?

Sva ova, ali i mnoga druga pitanja, mogu se postaviti kada govorimo o veštačkoj inteligenciji.

\subsection{Veštačka inteligencija: Definicija i značaj}

Da bismo uopšte mogli da potražimo odgovore na ova pitanja, potrebno je da prvo razumemo sa čime baratamo. Šta je, zapravo, \textbf{veštačka inteligencija}?

Veštačka inteligencija (AI) je tehnologija koja omogućava mašinama da obavljaju zadatke koji bi inače zahtevali ljudsku inteligenciju. To uključuje prepoznavanje govora, razumevanje slika, predviđanje ponašanja ili donošenje odluka. Snaga AI leži u njegovoj sposobnosti da prepoznaje obrasce i uči iz iskustava. Ovaj proces omogućavaju algoritmi koji se obučavaju pomoću velikih količina podataka~\cite{whitevision_ai_privacy}. Veštačku inteligenciju možemo shvatiti kao simulaciju ljudske inteligencije, a obučavanje algoritama možemo uporediti sa ljudskim učenjem.

AI brzo menja način na koji radimo, komuniciramo i obrađujemo informacije. Od pametnih algoritama koji preporučuju proizvode do \textit{chatbotova} koji pružaju korisničku podršku, AI postaje sve prisutniji u našim svakodnevnim životima~\cite{whitevision_ai_privacy}. Pored toga što nam AI donosi veliki broj novih mogućnosti, ubrzava rad i obradu informacija, pojavljuju se novi izazovi koje je potrebno prevazići. Jedan od njih, a možda i najvažniji, je \textit{privatnost}.


\subsection{Zašto je veštačkoj inteligenciji potrebno mnogo podataka}

Sistemi veštačke inteligencije zasnivaju se na velikim količinama podataka. \textbf{Veštačka inteligencija zavisi od podataka jer "`gradi svoje znanje"' analizom primera.}\cite{whitevision_ai_privacy}

Uzmimo, na primer, AI sistem koji razvrstava naučne tekstove prema oblasti. On se obučava "`čitanjem"' nekoliko hiljada naučnih tekstova. Koristeći te podatke, sistem prepoznaje obrasce koji određene tekstove grupišu u određene grupe. Pored toga, on pravi pretpostavke o tome koje su to bitne reči koje jedan tekst odvajaju od drugog. Što su podaci bolji, to će se više netačnih pretpostavki ispraviti ili odbaciti.\cite{whitevision_ai_privacy} Na primer, engleska reč "`The"' će sigurno biti prepoznata kao nebitna, dok će se, na primer, reč "`\textit{mikrokontroler}"' sigurno češće pojaviti u tekstovima vezanim za robotiku nego u tekstovima vezanim za medicinu.

Za preciznost i robusnost AI alata potrebne su ogromne količine podataka kako bi što bolje mogao da prepoznaje obrasce i postavi pretpostavke. Tu dolazimo do još pitanja. Odakle tolike količine podataka dolaze? Ukoliko se radi o ličnim podacima, da li imamo pravo da ih koristimo za treniranje naših alata?


\subsection{Uticaj veštačke inteligencije na privatnost}

Privatnost u eri veštačke inteligencije postala je jedno od najvažnijih pitanja u modernom društvu. Veštačka inteligencija (AI) se često povezuje sa problemima u vezi sa prikupljanjem podataka, sajber bezbednošću, dizajnom modela i upravljanjem podacima. Neki od ključnih rizika vezanih za privatnost u AI uključuju:

\begin{itemize}
    \item \textbf{Prikupljanje osetljivih podataka}.  
    Veštačka inteligencija donosi veće rizike po privatnost podataka nego prethodna tehnološka dostignuća zbog ogromnog obima podataka koji se koriste. Terabajti ili petabajti podataka, kao što su tekst, slike ili video zapisi, rutinski se koriste kao podaci za obuku. Ovi podaci često uključuju osetljive informacije, kao što su zdravstveni podaci, lični podaci sa društvenih mreža, finansijski podaci, biometrijski podaci koji se koriste za prepoznavanje lica i drugi. S obzirom na to da je količina osetljivih podataka koja se prikuplja, skladišti i prenosi veća nego ikad, rizik da će barem neki od tih podataka biti izloženi ili korišćeni na načine koji krše prava na privatnost takođe raste. \cite{ibm_ai_privacy}
    
    \item \textbf{Prikupljanje podataka bez pristanka}.  
    Rizici privatnosti mogu nastati kada podaci za razvoj AI tehnologija budu prikupljeni bez izričitog pristanka ili znanja ljudi od kojih se podaci prikupljaju. Na primer, korisnici internetskih platformi i sajtova sve više očekuju veću autonomiju nad sopstvenim podacima i veću transparentnost u vezi sa prikupljanjem podataka. Takve situacije su postale predmet spora, kao što je nedavni slučaj sa LinkedIn platformom, koja je suočena sa kritikama jer su neki korisnici primetili da su automatski uključeni u dozvolu za korišćenje svojih podataka za obuku generativnih AI modela. \cite{ibm_ai_privacy}
    
    \item \textbf{Korišćenje podataka bez dozvole}. 
    Čak i kada podaci budu prikupljeni uz saglasnost korisnika, rizici privatnosti ostaju ako se podaci koriste u svrhe koje nisu bile prvobitno obelodanjene. Na primer, fotografije ili biografski podaci koje smo podelili za jedan razlog, mogu se koristiti za obuku AI sistema, često bez našeg znanja ili pristanka. U nekim slučajevima, korisnici su otkrili da su njihovi medicinski podaci korišćeni za obuku AI sistema, iako nisu dali saglasnost za takvu upotrebu \cite{ibm_ai_privacy}.
    
    \item \textbf{Nekontrolisana nadglednost i pristrasnost}. 
    Problemi sa privatnošću koji se odnose na široko rasprostranjenu i nekontrolisanu nadglednost, bilo kroz sigurnosne kamere na javnim mestima ili kolačiće za praćenje na ličnim računarima, postojali su i pre širenja AI tehnologija. Međutim, AI može pogoršati ove probleme jer se AI modeli koriste za analizu podataka sa nadzora. U nekim slučajevima, rezultati takve analize mogu biti štetni, posebno kada pokazuju pristrasnost. Na primer, u oblasti sprovođenja zakona, brojna pogrešna hapšenja ljudi na osnovu AI odluka povezana su sa pristrasnostima u podacima. \cite{ibm_ai_privacy}
    
    \item \textbf{Ekstrakcija podataka}.
    AI modeli sadrže veliku količinu osetljivih podataka koji mogu postati plen za napadače. Na primer, u napadima na AI aplikacije, hakeri mogu koristiti tehnike kao što su prompt injekcija, gde maskirani ulazi manipulišu generativnim AI sistemima kako bi izložili osetljive podatke. Tako, hakeri mogu uputiti AI asistenta da podeli privatne dokumente, što može dovesti do curenja podataka. \cite{ibm_ai_privacy}
    
    \item \textbf{Curenje podataka}. 
    Curenje podataka je slučajno izlaganje osetljivih podataka. Neki AI modeli su se pokazali ranjivima na takve provale. Na primer, u jednom poznatom slučaju, ChatGPT je korisnicima pokazao naslove istorija razgovora drugih korisnika, što je dovelo do neželjenog deljenja privatnih podataka. Čak i mali, proprijetarni AI modeli mogu izazvati curenje podataka ako, na primer, AI aplikacija zasnovana na podacima korisnika nenamerno deli informacije sa drugim korisnicima koji koriste određeni upit. \cite{ibm_ai_privacy}
\end{itemize}

\section{Curenje podataka u sistemima veštačke inteligencije}

\textbf{Curenje podataka} (eng. \textit{data leakage}) predstavlja ozbiljan izazov u sistemima veštačke inteligencije. Osim što može negativno uticati na tačnost i performanse modela, može dovesti i do neovlašćenog otkrivanja privatnih ili poverljivih informacija, što otvara prostor za zloupotrebe.

\subsection{Uticaj curenja podataka na tačnost AI modela}

Iako ne utiče direktno na privatnost korisnika, \textbf{tačnost AI modela} je jedan od važnih faktora na koje utiče curenje podataka.

Alati veštačke inteligencije zahtevaju velike količine podataka kako bi njihovi rezultati bili što precizniji i pouzdaniji. Često se može desiti da se prilikom treniranja modelu daju podaci koji \textit{nije trebalo da budu uključeni}. Takav pristup može nehotice izazvati ozbiljne probleme, uključujući netačne procene performansi tokom eksperimentalne faze razvoja AI modela, kao i slabiju ponovljivost rezultata. \cite{apicella2024dontpushbuttonexploring} Sve ovo dovodi do toga da se kasnije smanjuje tačnost AI modela.

U ovom kontekstu, curenje podataka se često označava i kao \textit{curenje obrazaca} (eng. \textit{pattern leakage}) — situacija u kojoj se informacije koje bi trebalo da ostanu skrivene nehotice uvode u obučavajući skup podataka, čime se kompromituje objektivna procena performansi modela. \cite{apicella2024dontpushbuttonexploring}

Posledica curenja podataka tokom razvoja modela najčešće dovodi do dobijanja veoma optimističnih rezultata, koji se u kasnijoj primeni ne potvrđuju. Kako je kreiranje AI modela postalo sve dostupnije, sve češće se dešava da se u taj proces upuštaju osobe koje nemaju dovoljno znanja o njegovom funkcionisanju. Takvi korisnici, usled nedostatka znanja ili razumevanja, često preskaču ključne korake poput pravilne pripreme podataka ili selekcije karakteristika (eng. \textit{feature selection}), što povećava rizik od curenja i pogrešne evaluacije modela. \cite{apicella2024dontpushbuttonexploring}

Pored pogrešne evaluacije, često je za treniranje AI modela potrebno dosta vremena, a ukoliko se greška ne uoči na vreme, potrebno je ponovo prolaziti kroz dug i zahtevan proces treniranja.

Ipak, druga strana curenja podataka jeste ona u kojoj poverljivi ili privatni podaci dospevaju u skup podataka za treniranje modela, a kasnije su često dostupni i kroz sam model. Najčešće se u takvim slučajevima radi o generativnim modelima, koji na osnovu nekog korisničkog ulaza generišu izlaz u vidu \textit{teksta, slike, videa} i slično.

\subsection{Curenje podataka u generativnim alatima: ChatGPT ili ThreatGPT}

\textbf{Generativni alati veštačke inteligencije} (GenAI) poslednjih godina napravili su ogroman iskorak, značajno menjajući i ubrzavajući brojne poslovne procese. Iako su doneli mnoge koristi i unapredili rad u različitim oblastima, predstavljaju i potencijalnu pretnju kada ih koriste zlonamerne osobe.

Kada je reč o sajber bezbednosti, GenAI ima primenu i u \textbf{odbrani} i u \textbf{napadu}. Alati poput \textbf{ChatGPT-a} mogu pomoći stručnjacima u zaštiti sistema, analizom velikih količina podataka o sajber pretnjama – uključujući ranjivosti, obrasce napada i indikatore kompromitacije – što poboljšava sposobnost detekcije i odgovora. \cite{10198233} Ipak, isti alati mogu olakšati napadačima pripremu i izvođenje napada, čak i bez dubokog tehničkog znanja.

GenAI može efikasno analizirati \textit{log fajlove}, \textit{mrežni saobraćaj} i druge tehničke podatke tokom \textbf{sajber incidenata}, omogućavajući bržu i automatizovanu reakciju. Takođe se koristi u \textit{edukaciji korisnika}, \textit{podizanju svesti o bezbednosti} i \textit{pripremi na sofisticirane pretnje}. \cite{10198233}

U razvoju softvera, GenAI može doprineti pisanju bezbednijeg koda i generisanju test primera za otkrivanje ranjivosti. \textbf{Veliki jezički modeli} takođe mogu pomagati u definisanju etičkih smernica koje jačaju sajber odbranu. \cite{10198233}

Međutim, upotreba generativnih alata nosi i rizike. Prijavljeni su slučajevi u kojima je nepažljiva upotreba dovela do curenja poverljivih informacija. Iako modeli poput ChatGPT-a nemaju direktan pristup korisničkim podacima, postoji opasnost da se zaštitni mehanizmi zaobiđu.

Tokom obuke modela koriste se podaci sa različitih izvora – veb-sajtova, foruma, članaka – često bez saglasnosti autora. Procene ukazuju da skup podataka koji koristi ChatGPT premašuje 570~GB, što implicira obradu velike količine stvarnih korisničkih sadržaja. Ovo izaziva zabrinutost jer su lični komentari, blog postovi i recenzije korisnika potencijalno korišćeni bez eksplicitne dozvole. \cite{WU2024102}

Primer za to je partnerstvo kompanija \textit{Stack Overflow} i \textit{OpenAI} iz maja 2024. godine, kojim se predviđa korišćenje odgovora sa foruma za treniranje jezičkih modela. Ovo je izazvalo burne reakcije korisnika, od kojih su mnogi brisali svoje odgovore tvrdeći da nisu dali saglasnost. Sa druge strane, predstavnici Stack Overflow-a ističu da su odgovori kolektivno vlasništvo platforme. \cite{edwards_stackoverflow_2024}

Još jedan poznat slučaj dogodio se u kompaniji \textit{Samsung}, čiji su zaposleni koristili ChatGPT za pomoć u kodiranju, pri čemu su nenamerno podelili poverljive informacije. Ovi podaci su mogli postati deo modela i time biti izloženi riziku od javnog otkrivanja. \cite{10198233}

Takođe, curenje podataka pogodilo je kompaniju \textit{MediSecure} — australijsku službu za isporuku lekarskih recepata — u aprilu 2024. godine, pri čemu su poverljivi zdravstveni podaci završili na dark webu. \cite{nature_privacy_llm}

Posebnu zabrinutost izaziva mogućnost da prosečan korisnik, pažljivim formulacijama, izvuče poverljive informacije koje su greškom uključene u obuku modela. Kao reakciju na ovaj rizik, kompanije poput Samsunga razmatraju zabranu korišćenja alata poput ChatGPT-a u poslovnom okruženju. \cite{10198233}


\section{Manipulacija i lažiranje podataka}
Sve češće smo svedoci situacija u kojima je teško razlikovati stvarne i istinite informacije od lažnih. Bilo da je reč o lažnim vestima, fotografijama ili drugim sadržajima, tehnologije za manipulaciju podacima postale su izuzetno napredne. Razvoj veštačke inteligencije dodatno je ubrzao ovaj proces, omogućivši gotovo svakome da generiše fotografije sa licima stvarnih ljudi, kao i video-snimke u kojima se čini da određena osoba govori nešto što nikada nije izgovorila – i to samo na osnovu kratkog audio zapisa njenog glasa.

Ovakav tehnološki napredak otvara brojna etička i pravna pitanja, posebno u pogledu privatnosti, identiteta i poverenja u informacije koje svakodnevno konzumiramo.


\subsection{Lažiranje fotografija (Deepfakes)}

Generativni AI modeli koji rade sa slikama imaju veliki uticaj na privatnost — bilo da je reč o generisanju ili prepoznavanju lica. U nastavku su navedeni neki od primera koji to potvrđuju.

U januaru 2020. godine, novinarka koja se bavi privatnošću, \textit{Kashmir Hill}, objavila je članak u \textit{New York Times}-u, u kojem je opisala kompaniju \textbf{Clearview.AI} — firmu koja tvrdi da pomaže američkim organima reda u identifikaciji nepoznatih osoba, upoređujući njihove fotografije sa onima koje se nalaze na internetu. Kompanija koristi model za prepoznavanje lica treniran na milionima javno dostupnih slika sa mreže. \cite{lee2024deepfakesphrenologysurveillancemore}

Tokom 2021. godine, izveštaji su pokazali da su policijske uprave u brojnim američkim gradovima koristile \textbf{Clearview.AI} za identifikaciju pojedinaca, uključujući i učesnike protesta pokreta \textit{Black Lives Matter}. \cite{lee2024deepfakesphrenologysurveillancemore}

Godine 2022, umetnica iz Kalifornije otkrila je da su fotografije za koje je verovala da se nalaze isključivo u njenim privatnim medicinskim zapisima, bez njenog znanja i pristanka, završile u skupu podataka za treniranje \textbf{LAION}-a — baze podataka korišćene za obučavanje modela kao što su \textit{Stable Diffusion} i \textit{Google Imagen}. Umetnica, koja ima retko medicinsko stanje koje je želela da zadrži privatnim, izrazila je zabrinutost zbog potencijalne zloupotrebe generativnih AI tehnologija koje imaju pristup njenim fotografijama. \cite{lee2024deepfakesphrenologysurveillancemore}

U januaru 2023. godine, \textit{Twitch} strimerka \textit{QTCinderella} uputila je javni apel pratiocima na mreži \textit{X} (bivšem \textit{Twitteru}) da prestanu sa širenjem linkova ka nelegalnom sajtu koji je hostovao AI-generisanu deepfake pornografiju nje i drugih poznatih žena sa interneta. \textit{„Biti viđen 'nag' protiv svoje volje NE BI TREBALO DA BUDE DEO OVOG POSLA“}, izjavila je. \cite{lee2024deepfakesphrenologysurveillancemore}

Ovi, kao i mnogi drugi primeri, pokazuju koliko daleko upotreba alata veštačke inteligencije može ići i koliki uticaj može imati na privatnost pojedinaca.

Veliki napreci u veštačkoj inteligenciji u poslednjem periodu omogućeni su upravo zahvaljujući prikupljanju ličnih podataka, bilo to dozvoljeno ili nedozvoljeno. U radu \textit{"`Deepfakes, Phrenology, Surveillance, and More! A Taxonomy of AI Privacy Risks"'} \cite{lee2024deepfakesphrenologysurveillancemore} detaljno se opisuje da se za to često koriste sistemi video-nadzora. Korisnici često i nesvesno dozvoljavaju prikupljanje video i audio materijala sa svojih nadzornih kamera tako što koriste različite AI alate za prepoznavanje pokreta, procenu brzine automobila ili prepoznavanje neočekivanih događaja.

Jedna od aplikacija razvijenih specijalno za lažiranje ljudskih lica je \textbf{\textit{Faceswap}}. \textit{Faceswap} je deepfake aplikacija koja funkcioniše na operativnim sistemima \textit{Windows}, \textit{macOS} i \textit{Linux}. Prema njenoj \textit{GitHub} stranici, aplikaciju koristi najmanje hiljadu korisnika. Namenjena je širokoj upotrebi i dovoljno je jednostavna da je mogu koristiti i obični korisnici bez tehničkog predznanja. Imajući u vidu da se deepfake tehnologija može lako zloupotrebiti, ova dostupnost predstavlja ozbiljnu pretnju — jer veliki broj ljudi može da je koristi u neetičke svrhe. Deepfake danas postaje sve veći problem upravo zbog toga koliko uverljivo može izgledati. Ukoliko se ovakvi video-snimci budu masovno delili putem interneta, to bi moglo izazvati haos, jer su dovoljno ubedljivi da navedu ljude da poveruju u njihovu autentičnost. \cite{Nazar_2020}

\subsection{Lažiranje vesti}
Međutim, postoje i oni koji koriste veštačku inteligenciju za širenje \textit{lažnih vesti}. Pre nekoliko godina, društvene mreže bile su preplavljene lažnim sadržajem, što je podstaklo korisnike da postanu svesniji ovog problema i da preduzmu korake u borbi protiv njega. Danas postoji niz aplikacija koje, uz pomoć AI tehnologija, mogu da analiziraju da li je neka vest \textit{istinita} ili \textit{lažna}. \cite{Nazar_2020}

Ipak, situacija se promenila. Lažne vesti se više ne šire isključivo u tekstualnom formatu – sada se pojavljuju i kao video sadržaji, što otežava njihovu detekciju. Većina ljudi i dalje video zapis doživljava kao pouzdan dokaz istine. U prošlosti, montaža video snimaka bila je veština dostupna malom broju ljudi, ali sa pojavom Deepfake tehnologije, ovakav vid manipulacije postao je dostupan i široj javnosti. \cite{Nazar_2020}

Studija predstavljena u radu \textit{"`Artificial Intelligence and New Level of Fake News"'}\cite{Nazar_2020} uključuje upitnik sproveden među studentima. Svi ispitanici su naveli da je Deepfake tehnologija veoma rizična za zloupotrebu. Većina je takođe priznala da su makar jednom bili prevareni lažnim vestima na društvenim mrežama, ističući da Deepfake video može da zavara veliki broj ljudi i izazove društveni haos. Ipak, većina učesnika smatra da su u stanju da prepoznaju razliku između autentičnog i izmenjenog videa pažljivim posmatranjem i proverom informacija putem interneta. \cite{Nazar_2020}

Neki su dodatno izjavili da Deepfake snimci izgledaju toliko uverljivo da ostavljaju utisak da nisu izmenjeni ni na koji način. Pojedini učesnici su predložili zatvaranje pristupa softverima koji omogućavaju kreiranje Deepfake sadržaja, iako su svesni da takva mera nije dugoročno delotvorna jer \textit{ništa na internetu ne nestaje zauvek}. Drugi su naglasili značaj podizanja svesti o informacijama koje kruže internetom, čak i kada su praćene vizuelnim dokazima, dok je treća grupa preporučila edukaciju i širenje informisanosti o Deepfake tehnologiji kao ključnu strategiju za smanjenje rizika. \cite{Nazar_2020}

Iako Deepfake može da proizvede video zapise koje je teško razlikovati od stvarnih, postoje indikatori autentičnosti, kao što su \textit{artefakti} koji se pojavljuju u generisanim videima. Jedno istraživanje uspešno je detektovalo Deepfake sadržaje analizom upravo ovih vizuelnih nepravilnosti. Slično kao i kod prepoznavanja lažnih vesti, detekcija Deepfake sadržaja oslanja se na veštačku inteligenciju. Međutim, kako se ova tehnologija bude razvijala, biće neophodno da i alati za njeno prepoznavanje napreduju – u suprotnom, može se dogoditi da Deepfake sistemi koriste te iste podatke za sopstveno unapređenje, postajući još teži za otkrivanje. \cite{Nazar_2020}

Uz sve to, kontinuirano se razvija veliki broj različitih AI alata koji se mogu lako zloupotrebiti, što ukazuje na potrebu za pažljivim upravljanjem rizicima koje ovakve tehnologije nose. \cite{Nazar_2020}

\subsection{Lažiranje naučnih radova}

Jedna od oblasti koja izaziva zabrinutost jeste upotreba veštačke inteligencije za kreiranje \textit{lažnih naučnih radova} koji na prvi pogled izgledaju legitimno. Iako fenomen lažnih radova nije nov, razvoj AI tehnologija otvorio je nove mogućnosti za generisanje kvalitetnih lažnih radova u znatno kraćem vremenskom roku, pri čemu ih je sve teže otkriti. Ovo otvara važna pitanja o \textbf{\textit{integritetu naučnog istraživanja}} i \textbf{\textit{pouzdanosti objavljenih radova}}. \cite{doi/10.2196/46924}

Nekoliko studija je pokazalo potencijal veštačke inteligencije da generiše veoma uverljive \textit{lažne, nestručne članke}. Na primer, u jednom nedavnom eksperimentu, istraživači su koristili jezički model veštačke inteligencije pod nazivom \textit{GPT-2} (Generative Pre-trained Transformer 2) kako bi generisali lažan novinski članak koji je prihvaćen za objavljivanje u poznatom veb magazinu, a da nije otkriven kao prevara. Slično tome, u studiji koja je istraživala sposobnosti AI jezičkih modela u generisanju \textit{naučnih apstrakata}, istraživači su otkrili da su generisani apstrakti često bili neprepoznatljivi u odnosu na stvarne apstrakte i da su čak mogli da prevare i ljudske recenzente. \cite{doi/10.2196/46924}

U istraživanju predstavljenom u radu: \textit{"`Artificial Intelligence Can Generate Fraudulent but Authentic-Looking Scientific Medical Articles: Pandora’s Box Has Been Opened"'} \cite{doi/10.2196/46924}, korišćen je ChatGPT kako bi se generisao naučni tekst pod nazivom \textit{"`Effectiveness of deep brain stimulation for treatment-resistant depression: a randomized controlled trial"'}.  
Generisani tekst su nakon toga revidirali korišćenjem istog alata, a potom su ga testirali na alatima za proveru da li je tekst generisan uz pomoć veštačke inteligencije.  
AI detektor koji koristi \textit{Content at Scale} procenio je da je verovatnoća da je recenzija generisana pomoću ChatGPT-a \textbf{72\%}, što znači da je \textbf{"`veoma verovatno da je generisana veštačkom inteligencijom"'}. AI klasifikator teksta koji je razvio \textit{OpenAI} svrstao je recenziju generisanu ChatGPT-om u kategoriju \textbf{"`verovatno generisana od strane veštačke inteligencije"'}.

Nekoliko profesora iz različitih oblasti pregledalo je rezultujući rad.  
Jedan od najstarijih profesora neurohirurgije (DN) pregledao je AI-generisani članak i dao sledeće komentare: \cite{doi/10.2196/46924}

U celini, generisani članak pokazuje visok nivo tehničke stručnosti i autentičnosti. Ipak, uočeni su i određeni nedostaci i konkretne greške. Najizraženiji nedostatak jeste \textbf{kraća dužina članka} u poređenju sa sličnim radovima, kao i \textbf{ograničen broj referenci}. Za ovo je verovatno odgovorno ograničenje kontekstualne veličine modela, jer model može da obradi samo određenu količinu informacija u jednom trenutku.  
ChatGPT je značajno napredovao u poređenju sa ranijim modelima obrade prirodnog jezika (eng. \textit{natural language processing - NLP}) u razumevanju kontekstualnih odnosa između informacija koje se pojavljuju na različitim mestima u tekstu. Ovaj napredak se često pripisuje njegovoj sposobnosti da kompresuje prethodni kontekst i nadoveže nove informacije na njega. Međutim, i pored ovog napretka, model i dalje može imati poteškoća u obradi informacija koje ne mogu stati u njegovu latentnu reprezentaciju. \cite{doi/10.2196/46924}

Manji problem predstavlja nedostatak informacija o tome da li je studija registrovana na \textit{ClinicalTrials.gov}, kao i izostanak broja etičkog odobrenja. Još jedno ograničenje jeste to što trenutno dostupna verzija ChatGPT-a nije trenirana na podacima nakon septembra 2021. godine, te nije u mogućnosti da pruži informacije posle tog perioda (npr. novije reference). Prilikom provere citata i liste referenci otkrivene su značajne greške. Iako je \textbf{9 citata bilo tačno} u pogledu relevantnosti i unosa u literaturi, drugih \textbf{8 citata bilo je pogrešno}. \cite{doi/10.2196/46924}

Uvidevši napredak veštačke inteligencije u pisanju naučnih radova, možemo samo očekivati trenutak kada više neće biti moguće razlikovati pravi od generisanog naučnog rada.  
Tu se postavljaju nova \textbf{pravna i etička pitanja}:  
Ko preuzima kredibilitet za napisani naučni rad — da li \textbf{autor veštačke inteligencije} ili \textbf{osoba koja je rad generisala}?  
Da li se takav rad može prihvatiti kao \textbf{intelektualna svojina pojedinca}?  
Ova i mnoga druga pitanja sve će se češće javljati, ali je odgovore na njih veoma teško odrediti.

\section{Pravni okviri i regulative za zaštitu privatnosti}

Kada sagledamo sve navedeno, postavlja se veliki broj pitanja o tome kako zaštititi sopstvenu privatnost. Najidealnije bi bilo kada bi sva ta pitanja bila rešena na pravnom nivou, tako da korisnici, tokom upotrebe AI alata, ne bi morali da brinu o svojim pravima i podacima. Međutim, već samo definisanje, a zatim i sprovođenje potrebnih regulativa, predstavlja izuzetno složen izazov.

Osim zaštite korisnika, važno je razmotriti i kako zakonske regulative i pravni okviri utiču na sam razvoj AI alata. Jedna od ključnih regulativa na ovom polju jeste evropska Opšta uredba o zaštiti podataka \textbf{(General Data Protection Regulation – GDPR)}, koja postavlja visoke standarde u vezi sa obradom i zaštitom ličnih podataka građana Evropske unije.

\subsection{Da li nas zakon može u potpunosti zaštiti}

Smatra se da veštačka inteligencija (AI), slično kao i obrada \textbf{Velikih podataka (Big Data)}, predstavlja značajan izazov za primenu tradicionalnih principa zaštite podataka. Ovi sistemi mogu zahtevati razvoj novih rešenja u cilju očuvanja informacione privatnosti i zaštite drugih osnovnih prava. \cite{mitrou_gdpr_ai_2018}

Zaista, napredak u oblastima AI i robotike odvija se znatno brže od sposobnosti pravnog sistema da pruži odgovore na brojna \textit{etička}, \textit{pravna} i \textit{društvena pitanja}. \cite{mitrou_gdpr_ai_2018}


Opšta uredba o zaštiti podataka (GDPR) sadrži određene termine koji se odnose na internet, kao što su \textit{društvene mreže}, \textit{veb-sajtovi}, \textit{linkovi} i slično, ali ne sadrži termin „veštačka inteligencija“, niti izraze koji označavaju srodne koncepte poput inteligentnih sistema, autonomnih sistema, automatizovanog zaključivanja i izvođenja, mašinskog učenja, pa čak ni velikih podataka (eng. \textit{big data}).\cite{sartor_gdpr_ai_2020}

Ovo jasno ukazuje na to da GDPR uspešno adresira određena pitanja vezana za internet i digitalnu komunikaciju, ali se ne odnosi eksplicitno na izazove koje sa sobom donosi razvoj veštačke inteligencije. Ipak, mnoge odredbe GDPR-a su relevantne i za AI.

Međutim, ovo nije neuobičajeno za tehnologije koje se u velikoj meri oslanjaju na obradu podataka. Propisi poput Opšte uredbe o zaštiti podataka (GDPR) neizbežno zaostaju za tehnološkim razvojem, makar iz razloga što je izuzetno teško prilagođavati regulativu istim tempom kojim se razvija tehnologija. \cite{mitrou_gdpr_ai_2018}

Sa druge strane, reforma zakona radi usklađivanja sa savremenim tehnološkim inovacijama može biti kontraproduktivna. Takva nastojanja mogu dovesti do začaranog kruga, jer se tehnologija nastavlja razvijati i tokom samih procesa konsultacija i pregovora, čime se dodatno povećava rizik od pravne nesigurnosti. \cite{mitrou_gdpr_ai_2018}

Uzimajući sve u obzir, može se zaključiti da pravni okvir ne može u svakom trenutku pružiti potpunu zaštitu, jer će uvek postojati izvestan jaz između brzog tehnološkog razvoja i dinamike zakonskih regulativa.

\newpage
\section{Strategije za zaštitu privatnosti u eri veštačke inteligencije}

Iako postoji velika zabrinutost u vezi sa veštačkom inteligencijom, na kraju se postavlja pitanje: \textit{da li od nje ima više koristi ili štete?} Odgovor na to pitanje, ponovo, nije jednostavan. Na primer, možemo uporediti nastanak automobila sa pojavom veštačke inteligencije. Kada su automobili prvi put uvedeni, nisu postojala saobraćajna pravila koja bi usmeravala vozače i pešake, što je dovodilo do brojnih nezgoda. Međutim, ljudi su shvatili da automobili donose mnoge koristi i stvorili su saobraćajna pravila, čime su značajno povećali odnos koristi u odnosu na štetu.
Važno je napomenuti da, ukoliko želimo da se (gotovo) potpuno zaštitimo, određene korake moraju preduzeti i kreatori alata i sami korisnici.

Programeri modela imaju ključnu ulogu u uspostavljanju odgovornosti tokom celokupnog životnog ciklusa modela veštačke inteligencije. Njihove odluke postavljaju temelje za razvoj i kasnije funkcionisanje AI sistema – od izbora izvora podataka do oblikovanja arhitekture modela. Tokom procesa razvoja, programeri moraju dati prioritet etici, uključujući transparentnost, izbegavanje pristrasnosti, zaštitu privatnosti i uključivanje relevantnih aktera.
\cite{radanliev2023ethicsresponsibleaideployment}

Transparentnost je od suštinskog značaja za kreiranje modela sa interpretabilnim procesima donošenja odluka, u koje zainteresovane strane mogu imati poverenja. Programeri bi trebalo da budu u stanju da objasne kako je algoritam došao do određene odluke, čime se dodatno osiguravaju poverenje i odgovornost. \cite{radanliev2023ethicsresponsibleaideployment}

Programeri takođe moraju pravilno rukovati ličnim podacima, poštujući zakone o zaštiti podataka i primenjujući strategije poput diferencijalne privatnosti, kako bi zaštitili pojedinačne podatke, a da pritom očuvaju korisnost skupa podataka. Uključivanje različitih zainteresovanih strana, poput etičara i krajnjih korisnika, može pružiti bolji uvid u etička pitanja i najbolje prakse, čime se osigurava da model ispunjava potrebe i očekivanja svih aktera.\cite{radanliev2023ethicsresponsibleaideployment}

Sa druge strane, postoje metode koje korisnici alata veštačke inteligencije mogu primeniti, nezavisno od nivoa tehničkog znanja. Neke od takvih metoda su:

\begin{itemize}
\item deidentifikacija,
\item federativno učenje,
\item generisanje surogata,
\item diferencijalna privatnost,
\item kreiranje lokalnih modela.
\end{itemize}

Navedene metode predstavljaju često korišćene pristupe za zaštitu privatnosti prilikom korišćenja veštačke inteligencije u obradi elektronskih zdravstvenih zapisa. \cite{nature_privacy_llm}

\textbf{Deidentifikacija} predstavlja proces pripreme podataka kroz koji se uklanjaju ili sakrivaju lični podaci. Deidentifikacija je privukla veliku pažnju širom sveta, a sa napretkom tehnika velikih jezičkih modela, tačnost u pravilnom prepoznavanju i uklanjanju ličnih identifikatora uz očuvanje upotrebljivosti podataka za analizu i istraživanje značajno je povećana.\cite{nature_privacy_llm}

Neke studije predlažu upotrebu \textbf{federativnog učenja} (Federated Learning – FL) kako bi se poboljšale performanse modela uz očuvanje privatnosti podataka. FL je decentralizovani pristup koji omogućava više organizacija da zajednički treniraju AI modele bez razmene sirovih podataka. Umesto toga, model se šalje do lokalnih izvora podataka, gde se trenira lokalno, a zatim se rezultati (npr. parametri modela) vraćaju i agregiraju u centralni model. Na ovaj način, velike količine podataka iz različitih zdravstvenih ustanova mogu da doprinesu treniranju modela, dok se istovremeno zadržava lokalizacija podataka i umanjuju rizici po privatnost. Ova metoda je posebno značajna u kontekstu osetljivih podataka kao što su medicinski zapisi, jer omogućava naprednu analitiku i AI primenu bez ugrožavanja sigurnosti i poverljivosti podataka.\cite{nature_privacy_llm}

\textbf{Generisanje surogata} podrazumeva stvaranje zamenskih podataka koji oponašaju statističke karakteristike originalnih podataka, a da pri tome ne ugrožavaju privatnost pojedinaca. Ovaj proces zamene originalnih vrednosti surogatima omogućava da deidentifikovani podaci, bilo nestrukturirani ili strukturirani, zadrže integritet svoje strukture i značenja.\cite{nature_privacy_llm}

\textbf{Diferencijalna privatnost} se može definisati kao nemogućnost razlikovanja između dve raspodele verovatnoće izlaza. Na primer, zamislimo dve baze podataka koje sadrže informacije o zdravstvenom osiguranju, a koje se razlikuju za samo jednog pacijenta. Diferencijalna privatnost osigurava da, čak i uz pristup jednoj od baza, nije moguće sa sigurnošću utvrditi identitet pacijenta koji nedostaje. Ova neizvesnost se postiže dodavanjem kalibrisanog šuma u rezultate upita baze podataka, često pomoću Laplasovog mehanizma.\cite{nature_privacy_llm}

Studije su procenjivale izvodljivost korišćenja \textbf{lokalno implementiranih} laganih, otvorenih i oflajn velikih jezičkih modela u bezbednim okruženjima kako bi se prevazišli neki od problema privatnosti povezanih sa korišćenjem velikih jezičkih modela trećih strana. Ova strategija može se primeniti na strukturirane i nestrukturirane podatke. Lokalna primena velikih jezičkih modela može biti i računarski ograničena i zahtevna po resurse, posebno uz sve veći broj elektronskih zdravstvenih zapisa. Da bi se prevazišla ova ograničenja, jedan od pristupa jeste upotreba laganih lokalno implementiranih velikih jezičkih modela ili tradicionalnih modela zasnovanih na pravilima i mašinskom učenju za deidentifikaciju i generisanje surogata pre nego što se informacije proslede vlasničkim ili eksternim velikim jezičkim modelima.\cite{nature_privacy_llm}

Sve ove metode mogu značajno doprineti zaštiti privatnosti korisnika, posebno u komercijalnim sferama. Važno je da korisnici primenjuju određene taktike, bez obzira na to šta autor modela tvrdi o njegovoj sigurnosti, kako bi dodatno zaštitili svoje podatke.

\section{Zaključak i budući izazovi}

Bez obzira na sve, veštačka inteligencija značajno doprinosi ljudima – kako u poslovne, tako i u zabavne i edukativne svrhe. Iako postoje brojna pitanja vezana za privatnost, uz promišljenu upotrebu rizici se mogu svesti na minimum. Samim tim, ne treba se plašiti korišćenja alata veštačke inteligencije, već ih treba upotrebljavati obazrivo i smisleno.

Sa razvojem veštačke inteligencije dolaziće i novi izazovi. Mnogi poslovi biće automatizovani, a (skoro) svi podaci teći će kroz autonomne sisteme, što će otvoriti brojna nova pitanja koja u ovom trenutku ne možemo ni da zamislimo.
\newpage

\AtNextBibliography{\small}
\printbibliography


\end{document}